% !TeX program = pdflatex
% =============================================================================
%  UOS — Unified Oral Scene. Product white paper.
%  Single-file LaTeX source. Compiles on Overleaf with pdfLaTeX.
% =============================================================================
\documentclass[11pt,a4paper]{article}

\usepackage[utf8]{inputenc}
\usepackage[T1]{fontenc}
\usepackage{lmodern}
\usepackage[english]{babel}
\usepackage[margin=2.5cm]{geometry}
\usepackage{amsmath}
\usepackage{amssymb}   % \mathbb, usado en SE(3), SO(3) y R^3
\usepackage{microtype}
\usepackage{booktabs}
\usepackage{longtable}
\usepackage{array}
\usepackage{xcolor}
\usepackage{enumitem}
\usepackage{fancyhdr}
\usepackage{titlesec}
\usepackage{indentfirst}

\definecolor{uosblue}{HTML}{000000}
\definecolor{uosgrey}{HTML}{5A6570}
\definecolor{uosrule}{HTML}{C8D2D9}

\usepackage[colorlinks=true,linkcolor=uosblue,urlcolor=uosblue,citecolor=uosblue]{hyperref}

% Identificadores como `extensions_required` no se parten, y en una linea estrecha se
% salen al margen. Esto deja a TeX estirar el espaciado en la linea desesperada en vez
% de desbordarla: una linea algo suelta se lee, una que invade el margen no.
\setlength{\emergencystretch}{3em}

\setlist{itemsep=2pt,topsep=4pt,parsep=0pt}
\renewcommand{\arraystretch}{1.25}

% Encabezados de seccion sobrios, en serif y sin color: es un preprint, no un folleto.
\titleformat{\section}{\normalsize\bfseries}{\thesection.}{0.6em}{}
\titlespacing*{\section}{0pt}{1.6ex plus .6ex minus .2ex}{1.0ex plus .2ex}
\titleformat{\subsection}{\normalsize\itshape}{\thesubsection.}{0.6em}{}
\titlespacing*{\subsection}{0pt}{1.2ex plus .4ex minus .2ex}{0.7ex plus .2ex}

% Sin encabezado; el numero de pagina va CENTRADO AL PIE.
\pagestyle{fancy}\fancyhf{}
\fancyfoot[C]{\small\thepage}
\renewcommand{\headrulewidth}{0pt}
\renewcommand{\footrulewidth}{0pt}

\newcommand{\code}[1]{\texttt{\small #1}}
\newcolumntype{L}[1]{>{\raggedright\arraybackslash}p{#1}}
\newcolumntype{C}[1]{>{\centering\arraybackslash}p{#1}}

% a finding box: rules above and below, page-breakable
\newenvironment{finding}[1]
 {\par\medskip\noindent{\color{uosblue}\rule{\linewidth}{0.9pt}}\par\nobreak\vspace{3pt}
  \noindent{\bfseries\color{uosblue}#1}\par\nobreak\vspace{2pt}}
 {\par\vspace{2pt}\noindent{\color{uosrule}\rule{\linewidth}{0.6pt}}\par\medskip}


% =============================================================================
\begin{document}

\begin{center}
{\LARGE UOS: A Container for the Whole Dental Case}\\[0.30cm]
{\large an open format for cases that arrive in pieces}\\[0.9cm]
{\large Luis Garbayo Fernández\textsuperscript{a,b}}\\[0.5cm]
{\small\itshape In collaboration with \textsuperscript{a}ANFAIA \& \textsuperscript{b}Histora}\\[0.3cm]
{\small\today}
\end{center}

\vspace{0.6cm}
\hrule height 0.8pt
\vspace{0.35cm}

\noindent{\normalsize\bfseries Executive summary}
\vspace{0.15cm}

{\small
\noindent\textbf{The problem.} A complete dental case is today three or four unrelated
artefacts --- an intraoral scan, a cone-beam CT series, clinical photographs and written
reports --- each in its own coordinate system, each opened by different software. The
relations among them are not missing because nobody knows them. Somebody established every
one of them at the moment of acquisition. They are missing because nothing in the workflow
obliges anyone to write them down, so they live in a clinician's memory or inside a
vendor's product and \textbf{they do not travel with the data}. Every time the case changes
hands they are rebuilt from scratch, or not rebuilt at all.

\noindent\textbf{What we built.} UOS (Unified Oral Scene) is an open container format for
the whole case. It does not re-encode anything: acquired originals stay in their native
formats and are referenced by content hash, and what the container adds is the layer that
did not exist before --- which coordinate frame each asset lives in, what transform relates
them \emph{and how well it fits}, which visit each belongs to, what a clinician measured
against what a model inferred, and what any of it is worth. The specification, a published
JSON Schema, a validator and a conformance fixture bank are available under Apache~2.0.

\noindent\textbf{What it is worth today.} On a real clinical case the container regenerates
the intraoral scan to $4.59\times10^{-6}$\,mm against a 0.1\,mm budget, verifies a 397-slice
CBCT series file by file, and returns \emph{more} than went in: the same geometry with
per-vertex colour measured from the patient's own photographs, a tooth code on every vertex,
and a watertight arch a laboratory can print. The cross-modal alignment ships with its
residual, 0.666\,mm, and with the statement that no clinician has verified it --- which is
the difference between a number a clinician can act on and one they have to trust.

\noindent\textbf{The discipline behind every design decision.} Two things that are not the
same must never look the same. An automatic alignment nobody reviewed and one signed by a
clinician. A colour sampled from the patient and one from a fallback gradient. A tooth
boundary drawn by a network and one drawn by a person. Each distinction costs a field in a
manifest and buys a recipient the ability to decide what to trust, which is the only thing
that makes a case worth sending at all.

\noindent\textbf{What we are not claiming.} This is a proposal with a published schema and
a validator, not a standard, and it stays one until a second independent implementation
opens a container written by the first. Nothing produced by a model is presented as clinical
fact. The figures here come from one real case and three CBCT volumes; none of them is a
clinical study, and no claim is validated against patient outcomes. We report three lines
of work that did not reach their objective alongside the ones that did, because a format
whose purpose is to keep measurement and inference apart earns nothing if its authors
exempt themselves from the rule.

\vspace{0.25cm}
\noindent\emph{Keywords:} dental informatics, interoperability, digital twin, provenance,
CBCT, intraoral scanning, DICOM, glTF, FHIR, 3D Gaussian splatting
}

\vspace{0.35cm}
\hrule height 0.8pt
\vspace{0.7cm}

{\small\tableofcontents}
\newpage

% =============================================================================
\section{A case that arrives in pieces}\label{sec:problem}

\subsection{What a practice holds after one visit}

Consider what a dental practice actually holds after a single patient visit that involved
imaging. There is one or two STL files from the intraoral scanner, a directory of several
hundred DICOM slices from the cone-beam CT, a handful of clinical photographs in whatever
the camera wrote, and a PDF report. Four artefacts, four programs, four coordinate
systems, and no machine-readable statement of how any of them relate to any other.

The relations are not absent, somebody established them. Someone looked at the CBCT
beside the scan and knew which surface was which; someone read ``distal caries on 24'' and
knew which of the fourteen crowns on screen that named. But that knowledge lives in a
person or inside a vendor's software, and \textbf{it does not travel with the data}. When
the case goes to the laboratory, to a second opinion, or to a reviewer six months later,
the relations are reconstructed from scratch --- or, more often, not reconstructed at all,
and the recipient works with whichever artefact their software opens.

\subsection{What it costs, and who pays it}\label{sec:cost}

It is worth being concrete about the cost of the fragmentation, because ``the data is not
integrated'' is the kind of statement everyone agrees with and nobody acts on. In each of
the four cases below, the person who loses is not the person who could have prevented it.

\paragraph{The second reader.} The clinician who acquired the case has the relations in
their head and does not need them written down, which is exactly why they never get
written down. The cost lands on the second reader. A colleague opening the same case a week
later sees a grey arch in one program and a stack of slices in another, and must decide
whether a finding described in prose belongs to the crown they are looking at. In practice
they re-derive what they can and take the rest on trust.

\paragraph{The laboratory.} What reaches a dental laboratory is usually an STL and a
written instruction. The STL is geometry and nothing else: no colour, no units declared, no
statement of which tooth is which, no indication of what was measured versus what the
scanner interpolated across a gap. Everything the clinic knew that is not geometry has to
be re-expressed in prose, and prose does not survive a workflow with three handoffs.

\paragraph{The next visit.} Longitudinal questions --- has this margin receded, has this bone
level dropped --- are arithmetic if two visits share a coordinate frame and a research
project if they do not. Today they do not, so the comparison is either done by eye or
skipped. This is the loss we find hardest to justify, because both acquisitions already
exist and the only missing element is a transform with an error bar.

\paragraph{The archive.} A case archived as four loose artefacts is a case whose
interpretation depends on software that may not exist in five years. Formats outlive
vendors only when they are open and self-describing; a folder of files plus a convention
held in one company's product is neither.

\subsection{Why choosing a different format does not fix it}

The instinct is to look for an existing container. Each candidate fails on a different
axis, and the failures are instructive because they are not deficiencies: they are
consequences of what each format was designed for.

\begin{longtable}{L{3.0cm} L{12.2cm}}
\toprule
\textbf{Format} & \textbf{Why it does not carry a dental case}\\
\midrule
\endhead
\textbf{DICOM}~\cite{dicom} & Models acquisitions superbly and is the reason CBCT interoperates at all.
It has no primitive for a Gaussian field, no natural web transport, and its model of
``study'' does not extend to a surface scan plus a photograph plus a report as one
composed object.\\
\textbf{glTF}~\cite{gltf} & The right scene graph, mature tooling, an extension mechanism that works.
No volumes, no clinical metadata, no provenance, no notion of a regulatory boundary.\\
\textbf{OpenUSD}~\cite{usdz} & Composition is precisely what USD is for, and it solves layering better
than we do. It is foreign to the clinical ecosystem: no DICOM path, no FHIR path, and an
implementation burden a dental viewer cannot justify.\\
\textbf{FHIR}~\cite{fhir} & The right vocabulary for clinical facts and the right target for a
practice-management connector. \textbf{FHIR Release 4 has no resource for dental geometry}, and
\code{Media} is defined for photo, video and audio. Meshes end up as
\code{DocumentReference}, which is honest but is not interoperability.\\
\bottomrule
\end{longtable}

The observation that organises this document is that none of these formats is missing
\emph{data}. They are missing the \emph{relations}, and the relations are what a second
reader needs and what no format currently obliges anyone to state. That is why UOS is a
\emph{composition} format rather than another encoding, and why \S\ref{sec:related} is about
where UOS \emph{maps onto} these standards rather than where it replaces them.

\subsection{Why now}

Three things changed recently enough that the problem is worth attacking today rather than
five years ago.

\begin{enumerate}
  \item \textbf{Multimodal acquisition became ordinary.} A practice that had a CBCT and an
  intraoral scanner in different buildings a decade ago now has both, routinely, on the same
  patient in the same week. The relation between them stopped being an exotic research
  problem and became a daily one.

  \item \textbf{A representation appeared that can hold surfaces and volumes in one
  substrate.} 3D Gaussian splatting is fast, differentiable, and web-renderable, and it
  gives the project a common primitive for things that were previously meshes or voxels and
  never both. Whether that unification actually helps is an empirical question, and
  \S\ref{sec:gs} is our answer to it.

  \item \textbf{The regulatory boundary hardened.} As dental software incorporates models,
  the distinction between a measurement and an inference stops being an epistemic nicety and
  becomes the line that determines what may be distributed where. A format designed before
  that pressure has no place to put the distinction; one designed now has no excuse.
\end{enumerate}

\subsection{What this document is, and what it is not}

This is not the specification. The normative document is the \emph{UOS Format Specification
v0.3}~\cite{uosspec}, which defines the container, the manifest field by field, the
conformance levels and the validation algorithm, and which is accompanied by a published
JSON Schema, a reference implementation and a conformance fixture bank (\S\ref{sec:status}).
An engineer writing a reader should work from those.

What follows is the case for the format: the problem it exists to solve, the decisions
behind it and what they cost, and --- most of the document by weight --- \textbf{what has
actually been measured}, including what did not work. It is written for the two readers the
specification does not serve: the clinician or practice decision-maker who needs to know
what a case in this format is worth to them, and the partner or investor who needs to know
what exists, what is measured, and what is still a design rather than a result.

We follow one rule throughout, and state it here so the rest can be read against it.

\begin{finding}{Two things that are not the same must never look the same}
An automatic alignment nobody reviewed and one signed by a clinician. A colour sampled from
the patient's photographs and one from a fallback gradient. A tooth boundary drawn by a
network and one drawn by a person. A DICOM series held inside a container and one merely
identified by hash.

Each pair is easy to store identically and expensive to confuse. Where the format can
enforce the distinction it does; where it cannot, it is required to \emph{name} it. Most of
what looks like bookkeeping in \cite{uosspec} --- a residual on every registration, a
verification field defaulting to ``nobody'', a per-value declaration of how a clinical
figure was extracted --- is this one rule applied.
\end{finding}

% =============================================================================
\section{What UOS is}\label{sec:what}

\subsection{Composition, not encoding}

Our claim is that the missing artefact is a \emph{composition} format, and that composition
is a different problem from encoding. Encoding asks how to store a mesh or a volume well,
and it is largely solved. Composition asks what else must be written down so that the stored
things remain usable together: which coordinate frame each lives in, what transform relates
them and how well it fits, which visit each belongs to, which of them a model produced, and
what any of it is worth.

A \code{.uos} is therefore an uncompressed ZIP archive whose first entry is a
\code{manifest.json} --- the same trade \code{.usdz} makes \cite{usdz}, uncompressed so that
a viewer can read one layer out of a large case by HTTP range request instead of
downloading it whole --- and almost everything in it is metadata about payloads the format
does not re-encode. An acquired original --- the DICOM series, the scanner's own STL ---
\textbf{is referenced by its content hash and does not travel inside}. A directory asset
such as a CBCT series additionally carries a per-file digest list, so a recipient can prove,
slice by slice, that the series they hold is the one the case was built from, without the
container ever having held it. What \emph{does} travel is everything UOS composes: the
scene, the Gaussian layers, the clinical layer, the inference plane, the views, and the
relations that tie them together.

That choice has two consequences worth stating plainly. The container stays small enough to
send, and it never becomes a second copy of the patient's imaging archive with its own
drift. And what comes back out is \emph{not} the file that went in: it is the same geometry
regenerated from what the container carries, which is a stronger claim and a checkable one
(\S\ref{sec:evidence}).

\subsection{The three planes: measured, presented, inferred}\label{sec:planes}

Everything inside a container belongs to exactly one of three planes, and the separation is
the format's load-bearing structure rather than a filing convention.

\begin{longtable}{L{2.6cm} L{3.2cm} L{9.4cm}}
\toprule
\textbf{Plane} & \textbf{Where} & \textbf{What it is}\\
\midrule
\endhead
\textbf{Data} & \code{assets[]} & What was acquired, and what the pipeline computed from it
without a model. Immutable: a writer never edits an asset in place.\\
\textbf{Presentation} & \code{views.json} & What the clinician was looking at --- camera,
visible layers, clipping. Disposable and cheap to rewrite.\\
\textbf{Inference} & \code{derived/} & Anything a model produced. \textbf{Removable}:
deleting the directory and its manifest entries leaves a valid, complete case.\\
\bottomrule
\end{longtable}

Cutting across them, every value in a container declares one of three regulatory layers:
\textbf{1} acquired, or transcribed from a document a person signed; \textbf{2} computed by
a deterministic, reproducible procedure with no trained model in the path --- an automatic
registration, a format conversion, a colour measured off a photograph --- and obliged to
declare what it was computed from; \textbf{3} model output, which lives only under
\code{derived/}.

Layer 2 is not a softer layer 3, and defining it was not a bookkeeping exercise. Until it
existed, everything the pipeline computed shipped as though it had been \emph{measured}:
the ICP transform, the converted mesh, the subsampled density field. The distinction that
matters to a regulator is whether a trained model is in the path; the distinction that
matters to a reader is whether the value can be recomputed. Layer 2 answers yes to the
second and no to the first.

\begin{finding}{Why the third plane is what makes the format shippable}
Distributing the same case in a jurisdiction where the AI module is not cleared requires
deleting a directory, not re-exporting anything. That is what lets us say that layers 1 and
2 are \textbf{not Software as a Medical Device} and that only the layer-3 modules are.

A claim like that is worth nothing unless it is enforced, so the validator makes ``layer 3
outside \code{derived/}'' an error \emph{in both directions}: inference that escaped the
directory, and a directory entry that forgot to declare itself. It has caught us twice ---
a tooth code baked into the scene, and the same code riding inside two Gaussian payloads
--- which is the argument for checking rather than documenting.
\end{finding}

\subsection{Every relation carries its error and its author}\label{sec:relations}

The formal object a container stores is a labelled graph. Each asset declares exactly one
coordinate frame; each registration is an edge carrying a rigid transform between two
frames; and one frame, conventionally the intraoral scanner's, is canonical. Placing an
asset in the canonical frame means composing transforms along a path to it, which is well
defined only if such a path exists. That connectivity requirement is normative
\cite{uosspec}: an asset in an unreachable frame is not partially placed, it is one whose
position is undefined, and a viewer that draws it anyway is inventing a pose. The algebra is
in Appendix~\ref{app:math}.

What matters at this altitude is what rides on the edges. The CBCT-to-scanner edge on the
reference case was obtained by point-to-plane ICP over the extracted bone--enamel surface,
and the container ships the root-mean-square residual alongside it:
$\varepsilon_{\mathrm{RMS}} = 0.666$\,mm. That single scalar is what makes a cross-modal
measurement quotable. A crown-to-root length taken across the seam inherits it, and a
clinician deciding whether a 1\,mm difference is real needs to know that the alignment
itself moves by two thirds of that.

The container additionally records \emph{who} produced the estimate. On this case the
answer is a machine and nobody has verified it, so the registration is \textbf{provisional},
and both the validator and the viewer are required to present it as such. The same
registration also declares what it is fit for --- \code{visualization} --- because 0.666\,mm
is fine to look at and does not license surgical planning.

\begin{finding}{Why the graph is stored rather than collapsed}
It would be simpler to resolve every transform once and store assets pre-transformed. That
discards two things. First, the composition path: an error propagates along it, and a pose
obtained through two hops is not as trustworthy as one obtained through one. Second, and
decisively, the \emph{residual}. A collapsed geometry has no place to record that this edge
fits to 0.666\,mm and that one was drawn by hand.

The two failure modes are not symmetric. If the alignment is \emph{recomputed} by each
recipient, the case has as many geometries as it has readers, and two clinicians measuring
the same distance get different answers with no way to discover why. If instead it is
\emph{baked} into a merged artefact, there is one geometry but the residual is gone, so a
measurement across the seam carries an error nobody can quote. The second is the more
dangerous, because it looks like success.
\end{finding}

\subsection{What is inside one, in practice}\label{sec:inside}

The reference container used throughout this document holds a real maxillary case: an
intraoral scan, a CBCT, nine photographs and three reports, contributed by the collaborating
practice and processed pseudonymised. Opening it, this is what a reader meets.

\begin{longtable}{L{6.2cm} L{9.0cm}}
\toprule
\textbf{Item} & \textbf{What is there}\\
\midrule
\endhead
ZIP entries & 12 --- the acquired originals are referenced, everything UOS composes is
inside. Measured on the container as last written, before 0.6.0 gave the per-Gaussian tooth
labels their own \code{derived/} entries\\
Assets declared & 18, of which 13 are referenced rather than carried\\
Stored views & 18: four anatomical, the rest one per annotated tooth\\
Mesh & 112{,}067 vertices, 220{,}085 triangles, a single unpartitioned primitive\\
Appearance layer & 113{,}827 Gaussians, spherical harmonics degree 1,
\code{srgb\_rec709\_display}\\
CBCT density field & 1{,}341{,}990 Gaussians, subsampled $3\times3\times1$ from
12{,}072{,}785\\
Declared extensions & 3, all listed as \emph{used}, none \emph{required}\\
Provenance chain & 3 versions, each pointing at the hash of the one before\\
Conformance & UOS-Core + UOS-Vol, 0 errors\\
\bottomrule
\end{longtable}

Two of those rows carry more than they appear to. \textbf{None of the three extensions is
required}: everything we add adds information, and a conforming reader that understands none
of them must still be able to open the case --- otherwise an open format is one only its
issuer can read in full. And the container is \textbf{logically append-only}: modifying a
case is not editing it, it is writing a new version of the manifest that names the previous
one's hash. The chain is what makes ``this is the same case, later'' a verifiable statement
rather than a filename convention.

\subsection{How a new modality enters}\label{sec:newmodality}

Dentistry is not only surfaces and densities. Shade is colour under a stated illuminant;
periodontal status is a set of per-site measurements over time; and there is a growing body
of work on the oral microbiome, where the datum is neither a mesh nor a volume but a
per-site quantity attached to anatomy. A container that can only hold triangles and voxels
forces every such datum into a PDF, which is where clinical information goes to become
unqueryable.

``Extensible'' is easy to claim, so it deserves a concrete answer. Suppose a study wants to
attach microbial load per periodontal site to a case. Nothing about that is geometry, and no
existing dental format has a place for it. In UOS it requires four declarations and no
change to the format:

\begin{enumerate}
  \item An \textbf{asset} carrying the measurements, with its own hash, media type and the
  visit it belongs to.
  \item A \textbf{frame}, here the canonical one since sites are named anatomically
  rather than positioned, so that the connectivity rule is satisfied.
  \item An \textbf{extension} entry declaring the schema of what was added, listed as
  \emph{used} so that a reader without it still opens the case and can say
  what it left unread.
  \item A \textbf{descriptor} stating, per column, the unit, whether the value is measured
  or derived, and the vocabulary it uses, so that a consumer is never left inferring
  semantics from a column name.
\end{enumerate}

The same four steps admit shade measurements, periodontal charting, or a time series from
an intraoral sensor. What the format supplies is not a schema for each; it is the
obligation to declare one, and the guarantee that a reader who does not know it degrades
predictably instead of silently.

\begin{finding}{What this does not solve}
Declaring a modality is not validating it. UOS makes microbial load carryable, addressable
by anatomy, and traceable to whoever measured it. Whether that measurement means anything
clinically is outside a container format, and a format that implied otherwise would be
making exactly the category error it exists to prevent.
\end{finding}

% =============================================================================
\section{A case, end to end}\label{sec:endtoend}

The argument so far is structural. This section walks one real case through the whole path,
because the value of a composition format is not visible in any single stage of it --- it
is visible in what the sixth reader can do that the first could not.

\paragraph{1. Acquisition, unchanged.} The practice scans, takes the CBCT, shoots the
photographs and writes the report exactly as it does today. Nothing about UOS asks a
clinician to change how they work, and a format that did would not be adopted. The four
artefacts land as they always have: STL, a directory of 397 DICOM slices, nine JPEGs, three
PDFs.

\paragraph{2. The container is written.} The pipeline ingests the four modalities, aligns
the CBCT to the scan, labels the teeth, composes the scene and writes the \code{.uos}.
Ingestion of a complete set takes \textbf{12.7\,s} against a 60\,s budget; the whole run is
one command. The acquired originals stay where the clinic keeps them and are named inside
the container by their content hash --- the CBCT series with a hash per slice, all 397 of
them.

What the container now holds that did not exist before the run is the relational layer: the
CBCT-to-scanner transform with its 0.666\,mm residual and the statement that nobody has
verified it, the visit each asset belongs to, a tooth code on every vertex sitting in the
removable inference plane, and the report's findings anchored to the teeth they name.

\paragraph{3. It is validated, by anyone, without trusting the writer.} Validation is not a
courtesy the writer performs on itself. A reader runs the published validator over the file
and gets back a report with stable codes: whether the manifest matches its own schema,
whether every declared asset is present and hashes as promised, whether the frame graph is
connected, whether any inference escaped \code{derived/}, whether the provenance chain tells
the same story as the manifests. On the reference case: \textbf{0 errors}, conformance
UOS-Core + UOS-Vol.

The warnings are as informative as the errors, and the one this case raises is the right
one: the CBCT-to-scanner registration is automatic and unverified, so the validator names it
rather than leaving the obligation to whoever remembers to look.

\paragraph{4. A second reader opens it.} The viewer opens the container in the browser
without uploading anything. The arch appears with colour measured from the patient's own
photographs rather than the scanner's grey; the CBCT's Gaussian layers can be switched on
behind it; the views the first clinician stored --- four anatomical and one per annotated
tooth --- are named and reachable; selecting a tooth highlights it and raises the card of
what the report said about it.

The second reader is not re-deriving anything, which is the whole point. And the things
they should not simply trust are marked: the root under each crown is declared
\emph{reconstructed}, not measured; the alignment it rests on is declared provisional with
its residual in millimetres.

\paragraph{5. It goes to the laboratory, and more comes out than went in.} What the
container returns is not the file that went in. The same geometry regenerates to
$4.59\times10^{-6}$\,mm --- four orders of magnitude below the 0.1\,mm budget, and the
residual is the \code{float32} of the STL format itself rather than anything the round trip
introduced. On top of that geometry the laboratory gets per-vertex colour measured from the
patient, a tooth code on every vertex, and, where they want a physical model, an arch closed
as a watertight solid (12{,}025 closing faces, 11{,}936\,mm$^3$, 3.2\,s) plus one file per
tooth with a scanner-measured crown and a CBCT-reconstructed root.

That last file is where the composition earns its keep: a crown measured to tens of microns
by the scanner, joined to a root no scanner can see, from a volume no scanner can produce.
Neither modality yields it alone.

\paragraph{6. The patient comes back.} Six months later the practice scans again. Because
both visits declare their frames and the transforms between them, ``has this margin
receded'' becomes a subtraction between two versions of the same case rather than a new
study. The provenance chain says which version came first and proves neither has been
quietly edited, because a container is append-only and each manifest names the hash of the
one before it.

\begin{finding}{What the sixth step is worth, and what it still costs}
Step 6 is the one that cannot be bought today at any price, and it is also the one this
project has only half delivered. The mechanism is built and the chain verifies. But the
measurement a periodontist actually wants across the two modalities --- gingival margin to
alveolar crest --- does not yet work on real data, and \S\ref{sec:roadmap} reports exactly
how it fails and by how much. We separate the two because a format that made the mechanism
available is a real contribution, and claiming the clinical measurement on top of it would
be the kind of overreach the format exists to prevent.
\end{finding}

% =============================================================================
\section{What the recipient gets today}\label{sec:today}

\subsection{An improved scan back out}

An arch rendered from an intraoral scan is grey. The project measures colour per tooth from
the patient's own photographs --- the median of pixels per crown third, cervical, middle
and incisal, taken from the photograph that best sees each crown, expressed in CIELAB. On
the reference case \textbf{97.2\,\%} of vertices receive colour from their own tooth.

What the remaining 2.8\,\% does is more characteristic of the format than what the 97.2\,\%
does. Ninety-seven vertices inherit colour from the nearest measured neighbour and 3{,}041
are painted with a fallback gradient that is \emph{not the patient's colour}. That
distinction travels as a per-vertex flag, because otherwise ``I have no colour here'' and
``I have the wrong colour here'' are indistinguishable --- and the second one looks better.

\begin{finding}{Comparable between teeth, not on an absolute scale}
Flash falloff is removed using the patient's own gingiva as an internal reference, which
makes crowns \emph{comparable to each other}. The absolute level is not established,
because a clinical photograph carries no grey reference. The container therefore supports
``24 is darker than 23'' and refuses to support ``24 is an A3.5'', and says which is which
rather than shipping a number that looks like a shade match.

The route to the absolute scale is a capture-protocol change rather than an algorithmic
one: a grey card in frame fixes the illuminant and turns the same pipeline into a
calibrated measurement. This is the cheapest open item in the project.
\end{finding}

Alongside colour, the scan comes back with a tooth code on every vertex and, on request, as
a watertight solid a slicer will accept. The scanner does not produce any of the three.

\subsection{A cross-modal measurement with an error bar}

The measurements that justify a fused twin at all are the ones that exist in neither
modality alone. Gingival margin to incisal edge lives entirely inside the scan and is
available today at scanner accuracy, tens of microns, with no registration involved.
Cemento-enamel junction to alveolar crest lives entirely inside the CBCT. The third ---
\textbf{gingival margin to alveolar crest} --- has no 2D equivalent and cannot be taken
from either: the margin is soft tissue and does not appear on a radiograph, and the crest is
under the gingiva. It exists only in the composition.

What UOS contributes to that class of measurement is not the measurement. It is that any
quantity taken across the seam \emph{inherits a number}, the registration residual, and
that the number is in the file rather than in the head of whoever ran the alignment. A
1\,mm difference measured across an alignment that moves by 0.666\,mm is a different claim
from the same difference measured inside one scan, and today nothing in the workflow tells
the recipient which they are holding. \S\ref{sec:roadmap} reports where this measurement
stands, which is honestly: not yet working.

\subsection{A case that can be verified instead of trusted}

Every asset in a container is named by its content hash, and a directory asset such as a
CBCT series carries one hash per file. On the reference case that is \textbf{397 slices,
419 entries, 0 errors} --- and the checks run in both directions, so a slice that is present
but undeclared fails just as a declared one that is missing does. Both mean the series
coming out is not the series that went in.

This matters more than it sounds. A hash of the whole set answers ``this series does not add
up''. A hash per slice answers ``slice 214 does''. The first sends someone back to the
beginning; the second is actionable.

\subsection{A case that outlives its software}

A folder of four loose artefacts plus a convention held in one vendor's product is not an
archive. A container whose manifest is JSON against a published schema, whose payloads are
native formats with public specifications, and whose validator is open source, can be opened
in ten years by somebody with none of the original software. That is the weakest promise in
this section and the one with the longest horizon, and it is why the format is published
under Apache~2.0 rather than described in a brochure.

\subsection{And, for a research group, an auditable corpus}

A set of containers declares, per asset, what was measured and what was inferred, by which
model, at which weights hash. A dataset can therefore be filtered on those distinctions
without opening the payloads --- which is the difference between a corpus and a directory.
The failure this prevents is specific and silent: training the next model on the previous
model's output contaminates it in a way that is invisible in any metric computed on the
mixed corpus. Today no public dental dataset carries that separation.

% =============================================================================
\section{The representation choice: what Gaussian splatting adds, and what it does not}
\label{sec:gs}

A composition format has to decide what a composed scene is made of. The project adopted 3D
Gaussian splatting \cite{kerbl2023} --- fast, differentiable, web-renderable, and a single
substrate for things that were previously either meshes or voxels. We tested that appeal
rather than assuming it, and this section reports where it pays, because the answer
constrains how the format should be used and is easy to get wrong in a way that looks
impressive. The algebra behind both results is in Appendix~\ref{app:math}.

\subsection{Re-fitting one acquisition reallocates capacity; it does not create resolution}

The first question was whether splatting could recover detail the acquisition did not
capture --- the informal hope that an optimised continuous representation ``cleans up'' a
noisy surface or sharpens a radiograph. We simulated a $9\times9$\,mm patch of occlusal
surface containing three findings chosen so that each exists on a different physical
support: an occlusal fissure of 0.15\,mm (geometry), a white spot lesion of 0.60\,mm
(surface colour), and a dentinal lesion of 1.50\,mm under intact enamel (density only).
Each modality was simulated through its physical chain: Gaussian PSF, sampling at its own
pitch, then noise.

\begin{finding}{What this means for the format}
A finding that exists only as density is invisible to a surface scanner at any
reconstruction quality, and a fissure below the CBCT's sampling pitch is not recovered by
fitting Gaussians to the CBCT. What splatting contributes is the ability to \emph{allocate
capacity where the data has it}, and that pays exactly when the modalities being combined
have different resolving powers.

That condition holds in dentistry, where an intraoral scanner resolves surface detail an
order of magnitude finer than a CBCT resolves internal structure, and it is why the
representation is worth adopting here. The result is therefore not a limit on the method so
much as a statement of what the method is for: \textbf{composition across heterogeneous
sources}, which is also what the format is for. Any sharpening claim has to come from
adding a modality that carries the detail, not from re-fitting one that does not.

\textbf{The scope of this negative result, stated precisely.} Each modality here is
simulated as \emph{one} acquisition at \emph{one} sampling pitch, so what is measured is
that re-fitting a single acquisition does not recover detail below that pitch. It is not a
statement about the method in general. Multi-view reconstruction \emph{can} exceed the
resolution of any single view when the views overlap with sub-pixel displacements, which is
the mechanism of multi-view super-resolution. That regime is not measured here.
\end{finding}

\subsection{The capacity ceiling is in the renderer, and one thing moves it}

Fitting a single field to a real CBCT and varying the seeding budget from 20{,}000 to
418{,}727 primitives, reconstruction quality on held-out views does not improve: PSNR stays
between 40.2 and 43.2\,dB across the whole range, and beyond roughly 10{,}000 \emph{live}
primitives the optimiser prunes back to a similar working set regardless of what it was
given. Twenty times the budget buys nothing.

Appendix~\ref{app:math} shows why, and the reason is arithmetic rather than physical: the
usable range is $[\alpha_{\min}, \tau_{\max}] \approx [1/255,\, 9.21]$, fixed by the
rasteriser's numerics. Neither more budget, nor finer resolution, nor spatial concentration
widens it. What does widen it is splitting the volume into disjoint tissue layers, each
with its own density calibration, each spending that narrow window on a narrower band.

\begin{center}
\begin{tabular}{lrrrrr}
\toprule
\textbf{Case} & \textbf{1 field (dB)} & \textbf{5 layers (dB)} & \textbf{$\Delta$} &
\textbf{Live, 1 field} & \textbf{Live, 5 layers}\\
\midrule
A & 43.23 & 47.62 & $+4.39$ & 11\,013 & 176\,865\\
B & 42.02 & 44.10 & $+2.08$ & 22\,887 & 185\,787\\
C & 37.44 & 44.85 & $+7.41$ & 13\,672 & 133\,628\\
\bottomrule
\end{tabular}
\end{center}

With run-to-run noise of about 0.5\,dB the gain is signal, and it replicates on all three
volumes. It is not a parameter count: a dedicated field seeded at the \emph{same density}
as the global one, over the same region, scores 45.02\,dB against 25.56\,dB --- 19\,dB at
equal primitive count. Specialisation, not size, is what pays, and the mechanism is that
16 times more primitives survive the renderer's cull once each layer carries its own
calibration.

This also happens to be how a clinician would want to look at the twin --- enamel, dentine
and bone, separately --- so the decomposition is a product feature rather than a fitting
trick. And the remaining route is a product decision rather than a research one: a
rasteriser that integrates density directly, without the 8-bit cull and the \code{float32}
saturation that fix the ceiling, would make calibration a convenience rather than a
necessity. It is the single change that would most raise the ceiling for volumetric work.

\subsection{What this section commits the format to}

Two design decisions in \cite{uosspec} follow directly. The Gaussian layers a container
ships declare whether they are \textbf{measured} or \textbf{reconstructed}, and the
appearance layer is declared reconstructed --- because \S\ref{sec:negative} measures that it
is. And a layer declared measured must state the occupancy threshold that decided which
tissue appears in it, because otherwise a rendering choice masquerades as a finding.

Stated as a position: the honest use of splatting in dentistry is composition and capacity
allocation across modalities of differing resolution, not enhancement of a single
acquisition. A fitted field \emph{does} produce visually smoother output than the voxels it
was fitted to, and that smoothness is the representation's prior rather than recovered
signal. Under a metric that rewards plausibility it looks like enhancement; under one that
asks whether a 0.15\,mm fissure below the sampling pitch is now visible, it is not. Any
denoising claim requires a task-level evaluation against ground truth that this project has
not done, and we would not accept one from anybody else either.

% =============================================================================
\section{What has been measured}\label{sec:evidence}

Every figure below was measured on real clinical data, and each is stated with the data it
came from, because the project has several and they are not interchangeable. Unless a row
says otherwise, the source is \textbf{the reference case}: one maxillary arch with an
intraoral scan, a CBCT, nine photographs and three reports, contributed by the
collaborating practice and processed pseudonymised. The volumetric figures below come
from three CBCT volumes of a licensed cohort that cannot be redistributed. Figures from the public Teeth3DS+ corpus are named as such.

\subsection{The four commitments the project was set}

\begin{longtable}{L{7.2cm} L{2.4cm} L{2.4cm} L{2.4cm}}
\toprule
\textbf{Commitment} & \textbf{Target} & \textbf{Measured} & \\
\midrule
\endhead
Ingestion latency, complete set (STL + CBCT + report) & $<60$\,s & \textbf{12.7\,s} & met\\
Fidelity of the mesh regenerated from the twin & $<0.1$\,mm &
$\mathbf{4.59\times10^{-6}}$\,\textbf{mm} & met\\
Automated test coverage & $>80\,\%$ & \textbf{95.1\,\%} & met\\
Reliability of the ingestion agents & $>95\,\%$ & \textbf{93.8\,\%} ($N=16$) &
\textbf{not met}\\
\bottomrule
\end{longtable}

The fourth is published as it came out. The single failure is a clinical report scanned
\emph{without a text layer}: the agent extracts nothing and declares itself failed, which is
the correct behaviour. With $N = 16$ real cases one failure is 6.2 points, so the figure
sits below target by one case and fixing it is OCR rather than architecture. Both $N$ and
the percentage travel together on purpose --- a small $N$ that is declared is defensible;
hidden behind a percentage it is not. On the public Teeth3DS+ corpus the same mesh agent
runs 120 meshes at \textbf{100\,\%}.

\subsection{Composition and geometry}

\begin{longtable}{L{5.6cm} L{4.0cm} L{5.6cm}}
\toprule
\textbf{Measurement} & \textbf{Value} & \textbf{What it establishes}\\
\midrule
\endhead
CBCT $\to$ scanner registration & \code{icp\_surface},
$\varepsilon_{\mathrm{RMS}} = 0.666$\,mm over the matched correspondences, unverified &
That cross-modal alignment carries a usable residual, and that this one is provisional,
which the container states rather than hides.\\

Mesh reversibility & $4.59\times10^{-6}$\,mm against a 0.1\,mm budget &
That the scan is recoverable from what the container carries. The residual is the
\code{float32} of the STL itself, not a loss the conversion introduced.\\

Per-slice verification & 397 slices, 419 entries, 0 errors, 0 warnings &
That a DICOM series is verifiable file by file: a missing, altered or extra slice is
identified individually rather than as ``the series does not match''.\\

Watertight arch export & 12\,025 closing faces, 11\,936\,mm$^3$, 3.2\,s &
That the composed scene yields a printable solid the scan alone does not provide.\\

Container overhead of uncompressed storage & 183.2\,MB stored vs.\ 37.3\,MB deflated &
That random access by HTTP range is paid for, and how much: the \code{.ply} layers
compress to 13--15\,\%, the \code{.glb} only to 81\,\%.\\
\bottomrule
\end{longtable}

\subsection{Volumetric reconstruction}

\begin{longtable}{L{5.6cm} L{4.0cm} L{5.6cm}}
\toprule
\textbf{Measurement} & \textbf{Value} & \textbf{What it establishes}\\
\midrule
\endhead
Calibration of $\sigma$ & $19 \to 43$\,dB & That the dominant error was a units mismatch
against the rasteriser's window, not a modelling failure.\\

Additivity of disjoint layers & max.\ error $3.6\times10^{-7}$ on normalised $\tau$ &
That layer decompositions can be scored against the same image as the global field, which
is what makes the comparison honest.\\

Tissue decomposition & $+2.08$ to $+7.41$\,dB, three volumes & That specialisation lifts
the ceiling reproducibly.\\

Same-density dedicated field & 45.02\,dB vs.\ 25.56\,dB & That the gain is specialisation
and not primitive count: 19\,dB at equal seeding.\\

Live primitives after culling & $16\times$ more across five layers & The mechanism: each
layer's own calibration keeps primitives above $\alpha_{\min}$.\\
\bottomrule
\end{longtable}

\subsection{Appearance measured from clinical photographs}

Of 112{,}067 mesh vertices, \textbf{108{,}922 (97.2\,\%)} receive colour measured from the
photograph that best sees their own tooth; 97 inherit it from the nearest measured vertex,
and 3{,}041 are painted with a fallback gradient carried explicitly as a per-vertex
\code{measured} flag. Flash falloff is removed against the patient's own gingiva with
per-channel slopes of 0.649 / 0.577 / 0.677, which makes crowns comparable to each other
but does not establish an absolute scale (\S\ref{sec:today}).

\subsection{The gingival boundary, without a model}

Otsu thresholding on the $a^\star$ channel of the clinical photographs the container already
carries separates tooth from gingiva at \textbf{3.4--4.3\,$\sigma$} across four views, and
traces the gingival scallop tooth by tooth. No model, no training, no labels: a threshold.

We report it here rather than among the segmentation work because of what it establishes
jointly with \S\ref{sec:negative}: the reason geometric segmentation cannot find the crown
boundary is that \emph{the boundary is not geometric}, and the signal that does carry it is
already inside the container. That is an argument for composition made by measurement
rather than by design.

% =============================================================================
\section{What did not work, and what it cost the format}\label{sec:negative}

Three lines of work did not reach the outcome they were set up for. We report them with the
same weight as the rest, because a format whose purpose is to keep measurement and inference
apart earns nothing if its authors report only what worked, and because in each case the
measurement identified the next step more precisely than the success would have.

\subsection{Per-tooth segmentation on the intraoral scan}

\paragraph{What we measure.} The pipeline produces FDI codes for all 14 teeth of a maxillary
arch \emph{in the correct anatomical order}, and the boundaries are wrong. Eleven of the
fourteen crowns can be rejected on anatomy alone, comparing measured mesiodistal width
against anatomical tables, without needing ground truth at all.

\paragraph{Why, and what follows.} The boundary a crown ends at is not a geometric feature:
across the gingival margin the surface is locally continuous, so no purely geometric
criterion has the information to place it. That diagnosis is confirmed by the complementary
measurement in \S\ref{sec:evidence}: the same boundary is separable at 3.4--4.3\,$\sigma$ on
the colour channel of photographs the container already carries. The route forward is
therefore fusion rather than a better mesh segmenter: constrain the geometric label with
the photometric boundary, which is measured, unlabelled and already present. Anatomical
width tables give an independent rejection test that is available today and needs no model
at all.

\paragraph{What it cost the format.} Because tooth labels come from a model they belong to
the inference plane, while semantic picking in a foreign viewer is conventionally defined
over per-primitive metadata in the scene. For one version the container split the mesh into
one primitive per tooth so that any glTF viewer could select one --- and that put model
output inside a layer-1 asset, so deleting \code{derived/} no longer removed the inference.
It was reverted. The mesh now travels \textbf{unpartitioned}, the labels live only in
\code{derived/}, and picking is rebuilt by the reader from vertex index \cite{uosspec}.

The price is real and is accepted deliberately: a foreign glTF viewer opens the container,
draws the arch, and cannot select a tooth. Paying it is what makes the three-plane
separation of \S\ref{sec:planes} a property rather than a paragraph.

\subsection{A better classifier that made the real task worse}

\paragraph{What we measure.} The composite CBCT$+$scan pipeline delivers each tooth with its
root, and half the pieces come out too long: 27--37\,mm where a whole tooth measures 20--25.
The excess is alveolar bone attached below the apex. The hypothesis was precision, since the
first segmenter had recall 0.948 with moderate precision, so a more precise model should
trim the bone. The model that won the benchmark lost on the real task, and the defect that
motivated training it was unchanged.

\paragraph{Why, and what follows.} Voxel-level precision and apex localisation are different
objectives, and the benchmark rewards the first. Bone immediately below an apex is, to a
voxel classifier, indistinguishable from root: same modality, adjacent densities, no
boundary in the signal. The information that separates them is anatomical: roots
terminate, alveolar crest has a shape, and both have population statistics. The next step is
therefore a truncation criterion fitted on anatomy, evaluated on the quantity that actually
matters (tooth length against tables), rather than a further classification run scored on
voxel overlap. Choosing the evaluation metric to match the clinical quantity is the part
that was wrong, and it is fixable without any new data.

\subsection{Sharpness of the appearance field}

\paragraph{What we measure.} The appearance field reproduces its training views at
34--36\,dB PSNR and still looks soft. The project's sharpness measure is the
\emph{footprint}, the major axis of the median Gaussian in $\mu$m, with a criterion
of $<563\,\mu$m fixed \emph{before} running. Across seedings from 113\,000 to 1\,090\,000
primitives the criterion is not met, and more primitives do not move it: the median
footprint measures 845.9\,$\mu$m, 46.5\,\% of primitives sit more than 1\,mm from the
surface, and the field is anisotropic by $3.0\times$.

\paragraph{Why, and what follows.} The cause is measured rather than supposed, and it is
not a limit of Gaussian splatting. \textbf{The colour supervision is flat albedo}: the
training views are renders of the mesh painted with the 39 numbers of
\S\ref{sec:evidence} --- one measured median per crown per third, plus gingiva --- rather
than the photographs themselves, because
projecting the photographs onto the surface needs the camera pose the pipeline does not
compute (\S\ref{sec:evidence}). Against a piecewise-constant target the image loss cannot
distinguish a large diffuse disc from a small sharp one, so haze is the optimum of the
problem as posed, and the optimiser finds it \emph{whatever} the capacity, the strategy,
the render resolution or the added supervision. That hypothesis was falsified in three
formulations, including a run at 2048\,px with 1.09\,M primitives that reached 34.15\,dB
and still measured a 906\,$\mu$m footprint.

The consequence is that this result does \emph{not} measure what a photometric field can
resolve from real clinical views; it measures what one can resolve from a target that
carries no texture. The route that adds information is therefore to recover the camera pose
and supervise on the photographs' own pixels, which is a prerequisite (\S\ref{sec:evidence})
rather than a hyperparameter. Two cheaper routes reallocate existing information rather than
adding any: regularise the footprint directly, penalising the median major axis so the
optimiser trades reconstruction error against the quantity we actually care about; or seed
at a density matched to the view resolution rather than to the mesh.

\begin{finding}{Why these three are in the paper}
Each was expected to work, and reporting only the results that did would describe a pipeline
that does not exist. More practically, each is load-bearing for a design decision: the
segmentation boundary is why the format has a removable inference plane at all, and the
sharpness measurement is why appearance layers are declared \emph{reconstructed} rather than
\emph{measured}. A criterion fixed before running and not met is a result; the alternative,
renegotiating it afterwards, would have produced a paper with three more successes
and nothing to act on.
\end{finding}


% =============================================================================
\section{Regulatory posture}\label{sec:regulatory}

Nothing produced by a model is presented as clinical fact. Inference ships as
\code{investigational}, isolated in a removable directory, at a declared regulatory layer,
with the model and the weights that produced it named. This is not a compliance claim; it is
the minimum structure under which such a claim could later be made by somebody who has done
the clinical validation, and under which the same case can be distributed where the
inference module is not cleared.

The structure is what makes the claim checkable rather than asserted. Layers 1 and 2 ---
what was acquired, and what a deterministic procedure computed from it --- are \textbf{not
Software as a Medical Device}; only the layer-3 modules are. That statement holds only if
inference cannot leak out of \code{derived/}, so the validator enforces it in both
directions and has caught us doing exactly that twice (\S\ref{sec:planes}).

Alongside the conformance levels, which say whether a reader can \emph{open} a container,
the format defines an orthogonal profile that answers a different question: whether the
container is in a condition to \emph{leave} the organisation that issued it. It is the
question asked immediately before attaching a case to an email, and nothing in the levels
answered it. A container is distributable only if it declares how it was de-identified, what
purpose of use the patient consented to, keeps every layer-3 asset inside \code{derived/},
and validates against its own manifest. \textbf{Not being distributable is not an error}: it
is the normal state of a case inside the clinic. What matters is that the answer reaches a
person rather than being assumed.

% =============================================================================
\section{How UOS relates to what already exists}\label{sec:related}

The failure mode for a new format is to reinvent what a standard already defines and then
discover the incompatibility at integration. UOS is designed to sit \emph{between} the
standards rather than in place of them, and the specification names the destination for each
piece even where no connector exists yet --- naming it costs nothing now and saves an
incompatible extension later.

\begin{longtable}{L{4.6cm} L{10.8cm}}
\toprule
\textbf{What UOS carries} & \textbf{Where it maps}\\
\midrule
\endhead
The registrations & DICOM's \textbf{Spatial Registration Object}, which references two
Frame of Reference UIDs and carries the matrix, so a PACS can store, version and attach it
to the study. UOS stores the matrix row-major, which is how DICOM stores it too.\\
The referenced scanner mesh & DICOM defines \textbf{Encapsulated STL Storage}
(Supplement~205). When the recipient is a PACS, that is the destination rather than a loose
file beside it.\\
The provenance chain & FHIR \code{Provenance}, field for field: which entity, generated by
which agent, when, superseding which.\\
Which teeth a photograph shows & FHIR R5's \code{ImagingSelection}. Naming R5's destination
is what stops somebody inventing a private FHIR extension for something R5 already has.\\
The scene & A binary \textbf{glTF}, with the Gaussian layers hanging off it. Where Khronos
ratifies \code{KHR\_gaussian\_splatting} \cite{khrgs}, our fallback path for external layers
retires in its favour.\\
The clinical findings & Coded as \code{\{system, code, display\}} against SNOMED CT, with
tooth numbering to ISO 3950 \cite{iso3950}.\\
\bottomrule
\end{longtable}

\begin{finding}{Where a code is deliberately left empty}
The findings travel coded, and the SNOMED \code{code} is \code{null}. Assigning one is
clinical terminology work, and a guessed code is indistinguishable from a correct one to
the connector that resolves it. The same discipline applies to the FHIR mapping: UOS
declares the resource \emph{type} --- \code{Media} for the photographs, \code{ImagingStudy}
for the volume, \code{DocumentReference} for the mesh and for the case --- and leaves the
reference empty, because this case has not been through any practice-management system and
there is no identifier to cite. Inventing one would make a connector try to resolve it.
\end{finding}

The honest summary of the relationship is this. DICOM remains the right place for
acquisitions and is where the CBCT stays. FHIR remains the right vocabulary for clinical
facts and the right target for a practice-management connector. glTF remains the right scene
graph and UOS uses it directly. What none of them has, and what UOS adds, is the layer that
says how all of that relates --- and that layer is designed so that each piece of it can be
handed to the standard that will eventually own it.

% =============================================================================
\section{Status, and what openness means here}\label{sec:status}

A format is a claim about the future, so it is worth being exact about what exists today and
what is a statement of intent.

\begin{longtable}{L{4.6cm} L{10.8cm}}
\toprule
\textbf{Artefact} & \textbf{Status}\\
\midrule
\endhead
\textbf{The specification} & \emph{UOS Format Specification v0.3} \cite{uosspec}: a draft,
with the container, the data model, the manifest field by field, the conformance levels, the
validation algorithm and a cookbook of implementation recipes.\\
\textbf{The JSON Schema} & Published per version, one file for the manifest and one for the
validation report, with every field documented. The schema is generated from the same
contract the implementation uses, so the two cannot diverge silently.\\
\textbf{The validator} & Open source, 22 checks, each finding carrying a stable code
(\code{UOS-E-005}) so that a report can be cited and diffed rather than read.\\
\textbf{The conformance fixture} & Thirteen containers: two valid, ten with exactly one
defect each, one exercising forward compatibility --- with the expected verdict and the
reason stated beside each. Built by running the real writer over synthetic data, and
re-verified by the test suite on every run.\\
\textbf{The reference implementation} & A Python package that writes, reads and validates,
and the pipeline that produced every figure in this document.\\
\textbf{The licence} & Apache~2.0, covering the document and the reference implementation
alike.\\
\textbf{The media type} & \code{application/vnd.histora.uos}, \textbf{proposed and not
registered}.\\
\bottomrule
\end{longtable}

\paragraph{How the document changes, and who decides.} Changes are opened as issues and pull
requests against the public repository that hosts both this source and the implementation. A
change to the format and the change to the implementation that realises it belong in the same
request --- a proposal that does not say how it would be validated is not yet a proposal ---
and each must state the rule, where a validator would check it and at what severity, and
what it does to a container written under the previous version. The editor named on the cover
decides, in the open, on the record of the request.

\paragraph{And how a version is published.} The schema's identifier is pinned to a git tag
rather than to a branch, because an identifier that returns different content depending on
the day does not identify. That leaves a version under development needing an identifier that
resolves while its content is still changing, and a published one needing an identifier that
never changes --- two requirements that cannot share a name. So there are two tags: a
pre-release pointer, which is what the identifier names while the title page says
\emph{draft} and which may move, and the release tag, created once when the version is
published and never moved. Publishing is therefore an explicit act rather than a side effect
of pushing a tag, which is what it used to be, and once is what it costs.

\begin{finding}{What ``open'' does and does not buy, and why we say ``not a standard''}
This is a proposal with one implementation, not a standard with a committee, and calling it
either would be the same category error the format exists to prevent. Two things would change
that, in order. First, a \textbf{second, independent implementation} opening a container
written by the first --- until that happens, interoperability is asserted rather than tested,
which is why the fixture bank exists and why a distributable fixture built on openly
licensed data is on the roadmap. Second, a \textbf{registered media type} and a home for the
schema at a domain somebody holds.

What openness already buys is smaller and real: a recipient can validate a container without
trusting whoever wrote it, an integrator can implement a reader from a document rather than
from a support call, and a case archived today can be opened by somebody who has none of the
software that produced it.
\end{finding}

% =============================================================================
\section{Limitations}\label{sec:limitations}

\begin{itemize}
  \item \textbf{One clinical case.} The composition figures come from a single real case.
  The Gaussian-field results come from three CBCT volumes of a licensed cohort that cannot
  be redistributed. Neither is a clinical study, and no claim here is validated on
  patient outcomes. No clinical accuracy figure exists, because there is no ground truth for
  these patients: everything published is either an upper bound or a consistency check.
  \item \textbf{No distributable fixture of a real case.} The conformance fixture bank is
  synthetic. The measured container holds a real patient's case and is not published, so an
  implementer can check a reader against the fixtures but not against a clinical container.
  \item \textbf{Signals are unimplemented.} The asset class exists; no writer emits one, so
  the time-series path of \S\ref{sec:newmodality} is designed but untested.
  \item \textbf{Signatures are unspecified.} What is missing is not code but the decision
  of which key signs and where it lives; signing with an invented key would produce a
  container that appears signed.
  \item \textbf{The clinical act is not recorded.} The container says which software and
  which model produced a datum, and whether a human verified an alignment. It does not say
  who performed the study, under what protocol. That is a different axis from software
  provenance, and \S\ref{sec:roadmap} explains why it is a design decision rather than a
  field.
  \item \textbf{The media type is unregistered} and the format is not a standard. It is a
  proposal with a published schema and a validator, and it remains one until a second,
  independent implementation opens a container written by the first.
\end{itemize}

% =============================================================================
\section{Roadmap and outlook}\label{sec:roadmap}

What follows is not a list of features we would like to build. It is what clinicians asked
for once they saw the twin, plus what the measurements above say it would take to give it to
them. We separate the two deliberately: a request is a requirement, and whether it is
currently reachable is a measurement.

\subsection{The measurement clinicians actually want}

Validating whether a vertical gingiva-to-enamel dimension of the kind shown in commercial
2D demos is reproducible on a twin turned out to reveal that it is not one measurement but
\textbf{three}, and that they differ in what they need and in what they are worth.

\begin{longtable}{C{0.6cm} L{5.4cm} L{3.2cm} L{5.2cm}}
\toprule
& \textbf{Measurement} & \textbf{Needs} & \textbf{Status}\\
\midrule
\endhead
1 & Gingival margin $\to$ incisal edge & Scanner only & \textbf{Available.} Everything
happens inside a single mesh at scanner accuracy, tens of $\mu$m, with no registration
involved. This is the well-conditioned rung and it is implemented.\\
2 & Cemento-enamel junction $\to$ alveolar crest & CBCT only & Blocked on locating the CEJ;
see \S\ref{sec:roadmap}.2.\\
3 & Gingival margin $\to$ alveolar crest & Scanner \emph{and} CBCT & \textbf{The one with no
2D equivalent}, and therefore the one that justifies a fused twin at all. Blocked on
registration accuracy.\\
\bottomrule
\end{longtable}

Rung 3 is where the argument of this document meets clinical practice. It cannot be taken
from a radiograph, because the gingival margin is soft tissue and does not appear on one; and
it cannot be taken from a scan, because the alveolar crest is under the gingiva. It exists
only in the composition, which is precisely what the container is for.

\begin{finding}{And it does not yet work, in a way worth stating}
Computed today across the two modalities, the margin-to-crest distance comes out
\textbf{negative} on part of the arch: the crest is placed above the gingival margin, which
is anatomically impossible. A negative distance is a useful failure, because it is
self-evidently wrong rather than plausibly wrong, and it bounds the error: the registration
residual and the crest extraction together exceed the quantity being measured. Dedicated
per-tooth registration measures 0.45\,mm and its positive control does not yet close.

The immediate work is therefore not a new measurement but the two things it rests on:
tightening the alignment where the measurement is taken, and extracting the crest with a
criterion that is anatomical rather than density-thresholded. Neither needs new data.
\end{finding}

\subsection{What is measurably out of reach, and what replaces it}

Absolute gingival recession is defined against the cemento-enamel junction, and we measured
that the CEJ is not recoverable from either of our modalities. Not from CBCT by density:
restorative metal dominates the neighbourhood, and cervical enamel is thinner than the voxel.
Not from CBCT by shape: the detected point moves \textbf{6.6\,mm} depending on the threshold
chosen, which is not a measurement but a choice, and that route was measured and discarded.
Not from the intraoral scan: the cervical step lies below the triangulation noise of the
mesh, checked on an arch \emph{with} visible recession, so the failure is not an absence of
the feature.

What is solid is the gingival margin itself, detected in 11 of 11 sections across the three
arches tested. So the honest product is not absolute recession but \textbf{how far that
margin has moved between two moments}, which is also what a digital twin exists to provide:
not a photograph, a trajectory. Two visits in one frame make that a subtraction; what remains
is to establish the clinical threshold at which a displacement is reportable rather than
noise.

\subsection{Per-tooth motion, and a reference that does not move}

Orthodontic follow-up asks a narrower question: did \emph{this} tooth move, and by how much.
Per-tooth registration gives a robustly small residual, but the displacement it yields is
currently unusable, and the reason is instructive. Displacement is measured against the
global arch registration, and that registration itself shifts with the ICP trimming
parameter: the median displacement moves from \textbf{0.171 to 0.738\,mm} on that
hyperparameter alone. Every tooth appears to move because the frame moved. The reference has
to be built so that it cannot absorb the signal --- a leave-one-out frame, estimated from
every tooth except the one being measured.

One shortcut is worth recording as measurably wrong, because it is the kind of simplification
a roadmap adopts for performance reasons without checking: estimating the transform for three
teeth and averaging it for the rest is \emph{worse than doing nothing} in 2{,}629 of 2{,}860
comparisons, against simply using the global registration.

\subsection{Modalities that are not geometry}

\paragraph{The oral microbiome.} The clinical interest is not the bacteria as a curiosity but
which species are present and at what concentration, because several are associated with
conditions well outside the mouth and periodontal disease is the pathway through which that
association is thought to act. This is the concrete instance the extension mechanism of
\S\ref{sec:newmodality} was written for. Three requirements follow and only the first is
satisfied today: it must be addressable by anatomical site, which the container supports; it
must be temporally located, which the visit mechanism supports but no writer yet exercises;
and it must be declared \textbf{with its detection limit}, because a reported zero and a
concentration below the assay's threshold are different facts, and storing them identically
would repeat, on a new modality, exactly the confusion this project exists to avoid.

\paragraph{Occlusal force.} Clinicians also asked for bite force distributed over the
occlusal surface rather than as one number per patient. Mechanically this is the same shape
of problem as per-tooth colour: a scalar sampled at points and interpolated onto a surface the
container already carries. Two properties make it more delicate. A false-colour overlay must
never be mistakable for measured appearance, so a force map has to declare itself as a mapped
quantity with its own units rather than as an appearance layer. And the colour map is
\emph{part of the datum}: a field rendered through an unstated palette is not reproducible,
because the same measurements under a different scale tell a different clinical story.

\subsection{Audit: who did this, and how}

The container answers which software produced a datum, with which model and weights, and
whether a human verified an alignment. It does not answer what a clinician asks when a case
arrives from elsewhere: who performed this study, under what protocol, and when.

\begin{finding}{Why this cannot simply be a name field}
The obvious implementation, a professional's name on the acquisition, collides with the
property the format works hardest to preserve. A clinician is a person, their identity is
personal data, and a container that pseudonymises the patient while naming the practitioner
in clear has moved the problem rather than solved it.

So the work is a design decision with three parts: an actor identifier that can be resolved
by an institution and not by a recipient; a protocol reference that names how an acquisition
was performed rather than who performed it, since that is what a second reader needs in order
to reproduce it; and an extension of the declared PHI state to say whether professional
identities are present, so that the same honesty applied to the patient applies to everybody
in the record. Signatures are the natural carrier for the first, and are themselves
unspecified for the same reason: the missing piece is a decision about keys and custody, not
code.
\end{finding}

\subsection{Surgical planning, the furthest and the most valuable}

Graft planning is the case where the tissue decomposition of \S\ref{sec:gs} stops being a
reconstruction result and becomes the deliverable: what a surgeon needs to see is bone volume
available and bone volume required, separated from the soft tissue that hides it and from the
restorative metal that confounds it. The decomposition already produces those layers and the
container already carries them as separate, individually described fields. What is missing is
the clinical quantity computed on top of them, with its own error bar, and a validation that
it agrees with what a surgeon measures intraoperatively. We consider this the most valuable
open item and the one furthest from being closed.

\subsection{Platform work}

\begin{enumerate}
  \item \textbf{A rasteriser that integrates $\sigma$ directly}, without the 8-bit cull and
  the \code{float32} alpha saturation that fix the ceiling of \S\ref{sec:gs}. This is
  the single change that would most raise the ceiling for volumetric work, and it converts
  calibration from a necessity into a convenience.
  \item \textbf{A distributable conformance fixture of a realistic case}, built from an
  openly licensed dataset, so that interoperability can be tested rather than asserted.
  \item \textbf{Anatomical root truncation}, following \S\ref{sec:negative}: fit the
  truncation criterion on anatomy and score it on tooth length rather than voxel overlap.
  \item \textbf{Fusing the photometric boundary into the geometric segmentation}, using the
  colour separation of \S\ref{sec:evidence}, which is already measured and already inside the
  container.
  \item \textbf{A grey reference in the capture protocol}, which turns per-tooth colour from
  a comparative measurement into a calibrated one at no algorithmic cost.
  \item \textbf{A non-geometric modality end to end}, microbial load or periodontal
  charting, to test \S\ref{sec:newmodality} against a real dataset rather than a design.
  \item \textbf{Ratification tracking.} \code{KHR\_gaussian\_splatting} is a Khronos Release
  Candidate; when it ratifies, the container's fallback path for external Gaussian layers
  retires.
  \item \textbf{A second implementation}, ours or somebody else's, reading a container it did
  not write. This is the one that turns a proposal into a format.
\end{enumerate}

\begin{finding}{The pattern across all of these}
Every clinical item above is blocked on the same two things: the accuracy of a
cross-modality alignment, and the ability to locate an anatomical landmark that no single
modality resolves. Neither is a limitation of Gaussian splatting, and neither is solved by
more primitives. They are the reasons the format records a residual on every registration
and refuses to let an inferred landmark pass for a measured one, and they are where the next
year of work belongs.
\end{finding}

% =============================================================================
\section{Conclusion}

The problem in dental interoperability is not that the data cannot be encoded. It is that
the \emph{relations} between the encodings --- spatial, temporal, semantic --- and the
relation between a datum and the process that produced it are not written down anywhere that
travels. Every one of them exists at the moment of acquisition and is discarded before
delivery, because nothing obliges anyone to record it. We have described a container that
obliges it, and reported what building and measuring it taught us.

What exists today is a specification, a published schema, a validator, a conformance fixture
bank and a reference implementation, all open, plus a real clinical case that passes through
all of it: the scan regenerated four orders of magnitude inside its budget, a 397-slice CBCT
series verifiable slice by slice, an alignment that ships with its residual and with the
statement that nobody has verified it, and a scan that comes back out carrying more than it
went in with. What does not exist yet is a second implementation, and until it does this is a
proposal rather than a format.

The representational finding is the one we would most want carried forward, because it is the
one most likely to be claimed otherwise: Gaussian splatting is worth adopting in dentistry
for capacity allocation across modalities of different resolution, and not for enhancement of
a single acquisition, which our measurements show it does not deliver. The photometric
ceiling is real, budget does not move it, and specialised decomposition does, by up to
7.41\,dB at equal primitive count.

And the format's central commitment is a discipline rather than a feature: two things that
are not the same must never look the same. A registration nobody verified, a colour that is
the patient's versus one from a fallback gradient, a boundary drawn by a model versus one
drawn by a clinician. Each of those distinctions costs a field in a manifest and buys a
recipient the ability to decide what to trust, which is the only thing that makes a case
worth sending at all.

% =============================================================================
\newpage
\section*{Acknowledgements}

The author is grateful to ANFAIA and to HISTORA for their collaboration on this work.

The author thanks \textbf{Ismael Faro}, who founded ANFAIA and made this programme exist, for opening
the door to work of this kind in Galicia at all, and specifically for proposing that
Gaussian splatting be brought to bear on dental data in the first place. The idea that a
representation developed for novel-view synthesis might serve as a common substrate for a
surface scan and a volume is his, as is the practical framing that made it testable: fit it
to real acquisitions, on real hardware, and measure what it does rather than what it
promises.

Thanks are due also to \textbf{Matías Molinas}, CTO of
HISTORA, for accompanying the project throughout and for the technical judgement that
shaped several of the decisions defended in these pages, including more than one that
turned out to be a decision to measure something rather than to argue about it.

The author is indebted to the two clinicians who founded HISTORA. \textbf{Dr.\ Pedro Guitián}, oral
surgeon and founder of Centro Médico Odontológico Guitián, and \textbf{Dr.\ Elena López
Alvar}, oral surgeon, set out what a clinician actually needs to measure and, just as
usefully, what a plausible-looking number is worth when nobody can say how it was obtained.

Any errors, and every claim that turned out to be measurable rather than merely arguable,
remain the author's.


% =============================================================================
\appendix
\newpage

\section{Mathematical formulation}\label{app:math}

This appendix states the model the implementation rests on. It is here rather than in the
body because two of its results --- the capacity ceiling of \S\ref{sec:gs} and the
calibration that lifts it --- are consequences of the algebra rather than of any particular
dataset, and a reader who accepts the measurements does not need the derivation.

\subsection{The case as a graph}\label{app:graph}

The formal object UOS stores is a labelled graph. Let $F$ be the set of coordinate frames
named in a case and let one distinguished frame $f_0 \in F$, conventionally the intraoral
scanner's, be canonical. Every asset $a$ declares exactly one frame $\phi(a) \in F$.
Every registration is an edge carrying a rigid transform

\[
  T_{ij} \;=\;
  \begin{pmatrix} R_{ij} & t_{ij} \\ \mathbf{0}^{\!\top} & 1 \end{pmatrix}
  \;\in\; SE(3), \qquad R_{ij} \in SO(3),\; t_{ij} \in \mathbb{R}^3 ,
\]

mapping homogeneous points from frame $i$ into frame $j$, in millimetres, right-handed.
Placing an asset in the canonical frame is composition along any path to $f_0$,

\[
  T_{i \to 0} \;=\; T_{i_k 0} \, T_{i_{k-1} i_k} \cdots T_{i i_1},
\]

which is well defined only if such a path exists. That is the connectivity requirement the
format makes normative \cite{uosspec}.

\subsection{Registration as an estimator with a reported error}\label{app:icp}

The CBCT-to-scanner edge is obtained by point-to-plane ICP over the extracted bone--enamel
surface. Given source points $\{p_k\}$ and target points $\{q_k\}$ with target normals
$\{n_k\}$, the estimate minimises

\[
  \hat{T} \;=\; \arg\min_{T \in SE(3)} \;
  \sum_k \Big( \big( T p_k - q_k \big) \cdot n_k \Big)^{\!2},
\]

and the number the container ships alongside it is the root-mean-square residual

\[
  \varepsilon_{\mathrm{RMS}} \;=\; \sqrt{\tfrac{1}{K}\textstyle\sum_k \lVert \hat{T} p_k - q_k \rVert^2},
\]

measured on the reference case at $\varepsilon_{\mathrm{RMS}} = 0.666$\,mm.

\subsection{Identity of a directory asset}\label{app:hash}

Acquired originals are referenced by content address $H(\cdot) = \mathrm{SHA\text{-}256}$
and not carried. A directory asset such as a DICOM series additionally carries a per-file
digest list, and its own identifier is defined over names and hashes rather than
concatenated bytes,

\[
  H_{\text{series}} \;=\; H\!\Big( \big\Vert_{\,k \in \mathrm{sort}} \;
  \big( \mathrm{name}_k \,\Vert\, \texttt{NUL} \,\Vert\, H_k \,\Vert\, \texttt{LF} \big) \Big),
\]

so that renaming a slice changes the identifier, which is correct since slice ordering is
clinical data, and verification costs nothing beyond the manifest. The definition lives in
the contract rather than in the writer: if the validator computed it differently, a valid
container would come out invalid and nobody would know which of the two was right.

\subsection{Notation}
The model borrows from two literatures whose conventions collide on
one letter, and we keep both rather than renaming, so that a reader arriving from either
finds the symbol where they expect it. Lowercase $\sigma_i$ is the \emph{density} of
primitive $i$, the quantity that replaces opacity in the volumetric model. Uppercase
$\Sigma_i$ is its \emph{covariance matrix}, and uppercase $\Sigma$ without a subscript is
an ordinary summation. Standing alone, as in a ``$3\sigma$ cutoff'' or a separation of
``$3.4$--$4.3\,\sigma$'', $\sigma$ is a \emph{standard deviation} and has nothing to do
with density. Where a formula uses more than one of these, the surrounding text says which.

\begin{center}
\begin{tabular}{ll}
\toprule
\textbf{Symbol} & \textbf{Meaning}\\
\midrule
$c_i,\ \Sigma_i$ & Centre and covariance of Gaussian primitive $i$\\
$R_i,\ S_i$ & Its rotation and scale factors, $\Sigma_i = R_i S_i S_i^{\!\top} R_i^{\!\top}$\\
$\alpha_i$ & Opacity of primitive $i$, photometric model only, $\alpha_i \in [0,1]$\\
$\sigma_i$ & Density of primitive $i$, volumetric model only, $\sigma_i \ge 0$\\
$\mu(x)$ & Linear attenuation coefficient at $x$\\
$\tau(r)$ & Line integral of $\mu$ along ray $r$; $\tau_{\max}\approx 9.21$ is its
expressible ceiling\\
$k,\ k_\ell$ & Global and per-layer calibration scalars applied to $\sigma$\\
$T_{ij}$ & Rigid transform from frame $i$ to frame $j$, $T_{ij} \in SE(3)$\\
$\varepsilon_{\mathrm{RMS}}$ & Root-mean-square residual of a registration, in mm\\
\bottomrule
\end{tabular}
\end{center}

\subsection{The primitive}

A scene is a set of $N$ anisotropic Gaussians. Primitive $i$ has a centre
$c_i \in \mathbb{R}^3$ and a covariance factorised into rotation and scale so that it stays
positive semi-definite under gradient descent,

\[
  \Sigma_i \;=\; R_i S_i S_i^{\!\top} R_i^{\!\top},
  \qquad R_i \in SO(3), \quad S_i = \operatorname{diag}(s_{i1}, s_{i2}, s_{i3}),
\]

with $R_i$ parameterised by a unit quaternion and $s_{ij} > 0$. Its influence at a point is

\[
  G(x \mid c_i, \Sigma_i) \;=\;
  \exp\!\Big( -\tfrac{1}{2} (x - c_i)^{\!\top} \Sigma_i^{-1} (x - c_i) \Big).
\]

Rendering truncates at a Mahalanobis radius of $3$, the convention the Khronos extension
fixes \cite{khrgs}. This is the ``$3\sigma$'' of the standard-deviation sense above, not a
density: a primitive is treated as compactly supported on the ellipsoid
$(x-c_i)^\top \Sigma_i^{-1}(x-c_i) \le 9$.

\subsection{Two different physics on the same primitive}

\paragraph{Photometric splatting: appearance,} the original formulation \cite{kerbl2023}.
Each primitive carries an opacity $\alpha_i \in [0,1]$ and a view-dependent colour in
spherical harmonics. Pixels are formed by front-to-back alpha compositing over the
depth-sorted primitives covering them,

\[
  C \;=\; \sum_{i=1}^{N} c_i \, \alpha_i \prod_{j<i} \big(1 - \alpha_j\big),
\]

which is \emph{order-dependent}: the same set of primitives composited in a different order
gives a different pixel. The diffuse colour is recovered from the zeroth harmonic through
the basis constant $Y_0 = 1/(2\sqrt{\pi}) \approx 0.2820947918$ and the training bias,

\[
  \mathrm{RGB} \;=\; Y_0 \cdot \mathrm{SH}_{0,0} \;+\; \tfrac{1}{2}.
\]

\paragraph{Volumetric splatting: attenuation,} the clinical primitive. X-ray attenuation
is isotropic and unbounded, so opacity is replaced by a density $\sigma_i \ge 0$ and the
harmonics are discarded entirely. The field models the linear attenuation coefficient

\[
  \mu(x) \;=\; \sum_{i=1}^{N} \sigma_i \, G(x \mid c_i, \Sigma_i),
\]

and a detector value follows the Beer--Lambert law along the ray $r$,

\[
  \tau(r) \;=\; \int_r \mu(x)\, \mathrm{d}s,
  \qquad
  I \;=\; I_0 \, e^{-\tau(r)} .
\]

Because $\tau$ is an integral, this composition is \textbf{order-independent}: the same
primitives in any order give the same value, which is a materially different object from
alpha compositing, and it is why the container refuses to let the two share a vocabulary
without a descriptor saying which is which \cite{uosspec}.

\subsection{Where the capacity ceiling comes from}\label{sec:ceiling}

The clinical field is fitted with a rasteriser built for the photometric model, and that
borrowing has an algebraic price which we can state exactly.

Note that the field is expressed in $\sigma$ but rendered through a pipeline whose numerics
are defined on $\alpha$; the ceiling below is exactly the cost of that translation.
The rasteriser culls any primitive whose effective opacity falls below the 8-bit quantum,
$\alpha_{\min} \approx 1/255$. A culled primitive contributes nothing \emph{and receives no
gradient}: it is dead without notice. At the other end, accumulated alpha saturates in
\code{float32} at $\alpha_{\mathrm{acc}} = 0.9999$, which through
$\tau = -\ln(1 - \alpha_{\mathrm{acc}})$ caps the expressible line integral at

\[
  \tau_{\max} \;=\; -\ln\!\big(1 - 0.9999\big) \;\approx\; 9.21 .
\]

The usable dynamic range of the field is therefore the interval
$[\,\alpha_{\min},\, \tau_{\max}\,]$, and it is set by the renderer's numerics, not by the
anatomy. With $\sigma$ derived directly from Hounsfield units, hundreds of primitives per
ray saturate that ceiling: the model pins at \textbf{19\,dB} with predictions capped at
0.087 against a ground truth of 1.54.

\paragraph{The correction, and why it is legitimate.} A single global scalar $k>0$
rescales the field, $\sigma_i' = k\,\sigma_i$, calibrated so that the mean of the rendered
projection matches the mean of the digitally reconstructed radiograph. This is a choice of
\emph{units}, not of anatomy: it is a similarity on the range of $\mu$ and leaves the
\emph{shape} the CBCT measured untouched, so no structure is created or removed. With it,
reconstruction goes from \textbf{19\,dB to 43\,dB}.

\subsection{Additive decomposition into tissue layers}

Because $\tau$ is a line integral, it is additive over any partition of the field. If
$\{P_1, \dots, P_L\}$ partitions the primitives disjointly and exhaustively, then

\[
  \tau(r) \;=\; \sum_{\ell=1}^{L} \tau_\ell(r),
  \qquad
  \tau_\ell(r) \;=\; \int_r \sum_{i \in P_\ell} \sigma_i\, G(x \mid c_i, \Sigma_i) \,\mathrm{d}s ,
\]

so the layers must reproduce the complete radiograph exactly. We verified this before
comparing anything, measuring a maximum additivity error of $3.6 \times 10^{-7}$ on
normalised $\tau$, which is what makes it legitimate to score a single field and a
five-layer decomposition \emph{on the same image}, the only honest comparison.

The consequence is not merely bookkeeping. Each layer $\ell$ receives its own calibration
$k_\ell$, so a layer trained against the radiograph of one tissue attenuates far less and
survives the cull $\alpha_{\min}$ that killed it inside the global field. This is the
mechanism behind the gain reported in \S\ref{sec:gs}, and it is measured rather than
supposed: \textbf{16 times more live primitives} across five layers than in the single
field, at a seeding budget which, given to one field, does nothing.

\subsection{Appearance sampled onto a surface}

Recovering a per-vertex colour from a fitted field, the operation behind the improved
mesh, is a weighted mean over the primitives that actually cover a vertex $v$, weighted
as the rasteriser weights them. For the $k = 32$ nearest primitives within the $3\sigma$
support,

\[
  w_i(v) \;=\; \alpha_i \cdot G(v \mid c_i, \Sigma_i),
  \qquad
  \mathrm{RGB}(v) \;=\; \frac{\sum_i w_i(v)\, \mathrm{RGB}_i}{\sum_i w_i(v)} ,
\]

with a vertex marked unmeasured when $\sum_i w_i(v) = 0$. Using the nearest centre instead
would give the colour \emph{at} a point rather than the colour \emph{seen} at it, which
differs wherever primitives overlap: that is to say, everywhere.

\subsection{Thresholding the crown boundary}

Where a crown ends is ill-posed on geometry: a vertex on the gingival margin and one on the
cervical neck are locally the same surface. It is well posed on colour. Otsu's method
selects the threshold $t^\star$ on the CIELAB $a^\star$ channel maximising between-class
variance,

\[
  t^\star = \arg\max_t \; \omega_0(t)\,\omega_1(t)\,\big(\mu_0(t) - \mu_1(t)\big)^2 ,
\]

and we report the separation achieved as the standardised distance between the two class
means, where $\sigma_0$ and $\sigma_1$ are the within-class standard deviations and not
densities, $d = |\mu_0 - \mu_1| / \sqrt{(\sigma_0^2 + \sigma_1^2)/2}$, measured at
$3.4$--$4.3\,\sigma$ on real clinical photographs. No model, no training, no labels.

\subsection{Reporting metric}

Reconstruction quality is peak signal-to-noise ratio on held-out views,
$\mathrm{PSNR} = 10 \log_{10} (M^2 / \mathrm{MSE})$ with $M$ the dynamic range. Every
figure below is computed on views excluded from fitting (252 views per case, one in
eight retained), and every comparison between models is made on the same image.

% =============================================================================

\subsection{The seeding sweep, measured}\label{app:sweep}

The empirical counterpart of \S\ref{sec:ceiling}. A single field fitted to one real CBCT,
varying only the seeding budget, PSNR on held-out views:

\begin{center}
\begin{tabular}{rrrr}
\toprule
\textbf{Seeded primitives} & \textbf{Occupancy covered} & \textbf{Live primitives} &
\textbf{PSNR (dB)}\\
\midrule
20\,000  & 61.6\,\% & 7\,685  & 43.07\\
40\,000  & 76.6\,\% & 9\,373  & 43.16\\
80\,000  & 86.2\,\% & 11\,078 & 43.13\\
160\,000 & 93.3\,\% & 10\,721 & 40.55\\
320\,000 & 97.4\,\% & 8\,287  & 40.18\\
418\,727 & 98.0\,\% & 8\,298  & 42.44\\
\bottomrule
\end{tabular}
\end{center}

Twenty times the budget does not improve reconstruction, and beyond roughly 10\,000 live
primitives the optimiser prunes back to a similar working set regardless of what it was
given. The bound is the usable range $[\alpha_{\min}, \tau_{\max}]$ derived above, which is
fixed by the rasteriser's numerics and is widened by neither budget, nor resolution, nor
spatial concentration --- only by splitting the range across specialised fields
(\S\ref{sec:gs}).

% =============================================================================
\begin{thebibliography}{9}

\bibitem{uosspec}
L. Garbayo Fernández.
\emph{UOS: Unified Oral Scene: Format Specification v0.3}, 2026.
The normative document for the container described here: the data model, the manifest,
scenes, volume, the clinical and inference planes, provenance, extensions, the conformance
levels, the validation algorithm and implementation recipes. Published with a per-version
JSON Schema and a conformance fixture bank under Apache~2.0.

\bibitem{kerbl2023}
B. Kerbl, G. Kopanas, T. Leimkühler and G. Drettakis.
\emph{3D Gaussian Splatting for Real-Time Radiance Field Rendering}.
ACM Transactions on Graphics, 42(4), SIGGRAPH 2023.

\bibitem{gltf}
Khronos Group. \emph{glTF 2.0 Specification}.
\url{https://registry.khronos.org/glTF/specs/2.0/glTF-2.0.html}

\bibitem{khrgs}
Khronos Group. \emph{\code{KHR\_gaussian\_splatting}}, Release Candidate, 2026.
\url{https://github.com/KhronosGroup/glTF/tree/main/extensions/2.0/Khronos/KHR_gaussian_splatting}

\bibitem{dicom}
NEMA. \emph{DICOM Standard}, PS3.10 (media storage), PS3.15 Annex E.1
(attribute confidentiality profiles), Supplement 205 (Encapsulated STL Storage).

\bibitem{fhir}
HL7. \emph{FHIR Release 4} (4.0.1), 2019, and \emph{Release 5} (5.0.0), 2023.

\bibitem{iso3950}
ISO 3950:2016. \emph{Dentistry --- Designation system for teeth and areas of the oral
cavity.}

\bibitem{usdz}
Apple. \emph{USDZ File Format Specification}.
\url{https://openusd.org/release/spec_usdz.html}

\end{thebibliography}

\end{document}
